\subsection{Balance-Law Discretization and Implementation Verification}
\label{subsec:verification-literature}

One-dimensional arterial solvers are balance-law solvers: fluxes, sources,
boundary conditions, and time integration all affect the reported solution.
Finite-volume methods are common because they expose conservation and use
numerical fluxes to control wave propagation, while benchmark studies also
compare finite-element, DG, finite-difference, and network-oriented schemes
\cite{LeVeque2002FiniteVolumeMethodsForHyperbolicProblems,BoileauEtAl2015Benchmark,Wang2014BloodFlowNetworks}. Open implementations and recent
Lax--Friedrichs-type workflows provide useful comparators for reproducible solver
records
\cite{BenemeritoEtAl2024OpenBF,BeckersKolbe2025LaxFriedrichsHemodynamics}.

Well-balanced methods address a distinct requirement: a meaningful equilibrium
should remain steady under the discrete operator. This matters for stenotic
area--flow equations because geometry and wall-source terms can cancel in the
continuous balance while failing to cancel after discretization. High-order and
well-balanced methods therefore motivate explicit equilibrium tests, source
quadrature audits, boundary-regime diagnostics, and positivity or limiter records
\cite{PimentelGarciaEtAl2025HighOrderWellBalanced}. The present report uses that
literature to interpret the non-well-balanced geometry-rest drift as a material
numerical limitation.

Implementation verification asks whether the code solves the declared equations
and produces reproducible numerical records. Relevant evidence includes MMS
studies, exact-equilibrium tests, grid and timestep sensitivity, CFL and
subcriticality diagnostics, conservation records, positivity handling, and
artifact provenance. A benchmark-stage completion row or backend-parity check can
support implementation health, but it does not replace a manufactured-solution
record or a rest-equilibrium preservation test for the selected stenosis solver.

\subsection{Cross-Dimensional Observation Operators and Evidence Categories}
\label{subsec:cross-dimensional-comparison-review}

A 1D--3D comparison is meaningful only after both outputs have been mapped to a
common observable. For the velocity comparison in this report, the declared
plane--tetrahedron operator reduces the resolved 3D field to section area,
physical flow, and area-mean axial velocity, while the 1D state supplies
$Q_{1D}=\pi q$ and $\bar u_{1D}=q/a$. Agreement or discrepancy under this
operator is not automatically agreement under another observable.

The literature categories also need to stay separate. Implementation verification
checks a code against its declared equations and reference tests. Diagnostic
cross-model comparison reports discrepancies between two computational outputs
under a declared operator. External validation would require an accepted target
for the question, including matched geometry, wall and material assumptions,
boundary histories, timing, and measurement or reference conventions. The C23/C40
comparison in this report belongs to the second category.
